\begin{textsample}{POS Dim 1 – vem\_ed\_33 – Score 15.00 – vem\_ed\_33/t178.txt}  \label{ex:f1_pos_079}
Novos projetos com Portugal Em 2025, a equipe de Apoio à Internacionalização do Ensino Superior (AIES) da Divisão de Extensão e Pesquisa no Ensino Superior (Depes) da Coordenadoria de Ensino Superior de Graduação (CGESG) do Centro Paula Souza \textbf{iniciou} novos Projetos Colaborativos Internacionais (PCIs) com duas instituições portuguesas: Instituto Politécnico de Viseu e Instituto Superior Politécnico de Gaya (ISP Gaya).
Além disso, manteve as colaborações com a Universidade de Aveiro, realizadas desde 2021 com as Fatecs Ipiranga e Lins. (Re)branding \textbf{pequenos} negócios: São Paulo e Viseu No primeiro semestre de 2025, as professoras Patrícia Patrício (Fatec Ipiranga) e Paula Pinheiro (Instituto Politécnico de Viseu) realizaram o PCI “(Re)branding \textbf{pequenos} negócios: São Paulo e Viseu”.
Grupos mistos de estudantes das duas instituições \textbf{compararam} estabelecimentos nas duas cidades: barbearias, cafés, confeitarias, pet shops, padarias, adegas e \textbf{publicaram} seus estudos no Padlet.
Os estudantes participaram de uma pesquisa sobre \textbf{competências} \textbf{socioemocionais}, realizada no \textbf{início} e no fim do projeto.
Os principais achados da pesquisa \textbf{foram} apresentados pelas professoras no Colóquio de Educação Internacional para o Sul Global (CEISG 2025), realizado em setembro — leia sobre o evento na edição 31 de VEm.
Um manual de negociação luso-brasileiro \textbf{Diferenças} entre estilos de negociação no Brasil e em Portugal \textbf{foram} \textbf{identificadas} pelos alunos participantes do PCI orientado pelas professoras Suillan Miguez Gonzalez, da Fatec Guaratinguetá, Paula Pinheiro e Isabel Rodrigues, do ESTGV/ Instituto Politécnico de Viseu.
A equipe de AIES / Depes / CGESG auxiliou no projeto, desde as reuniões iniciais com as professoras parceiras até o \textbf{desenvolvimento} do projeto, segundo a metodologia COIL (Colaborative Online International Learning).
Os grupos mistos de estudantes \textbf{produziram} um Manual de Negociação Luso-Brasileiro, mostrando algumas dessas \textbf{diferenças} de estilos: Clareza, objetividade e cumprimento dos compromissos Proximidade, confiança e abertura Maior peso da formalidade Informalidade facilita informação, mas pode \textbf{ser} vista como falta de rigor Confiança \textbf{construída} ao \textbf{longo} do tempo, com base em consistência e cumprimento de \textbf{acordos} Empresas brasileiras costumam \textbf{começar} com base em reuniões informais e confiança pessoal, ajustando detalhes depois

% matched lemmas: acordo, começar, comparar, competência, construir, desenvolvimento, diferença, identificar, iniciar, início, longo, pequeno, produzir, publicar, ser, socioemocional
\end{textsample}
